\subsection{The exact character calculation for adding a puncture}

\begin{statement}{A.57}{Proved}{Joint fibre characters and their two-primary bound}
Let $4\mid n$, let $D_1,\ldots,D_k\in\Q^\times$,
and let
\[
 \mathcal E=\Q[Z_1^{\pm1},\ldots,Z_k^{\pm1}]
                    /(Z_i^n-D_i:1\le i\le k).
 \tag{A.57.1}
\]
This is a finite etale algebra. The finite-order
abelian characters of $G_\Q$ that vanish on the
absolute Galois group of every field factor of
$\mathcal E$ have the following exact finite description.

Write
\[
 D_i=(-1)^{\delta_i}\prod_{P\in\mathcal P}P^{w_{iP}},
 \quad \delta_i\in\{0,1\},\qquad
 F=\Q(\mu_{2n},\sqrt P:P\in\mathcal P),\qquad
 \Gamma=\operatorname{Gal}(F/\Q).
 \tag{A.57.2}
\]
For $\sigma\in\Gamma$ define $u_\sigma\in(\Z/2n)^\times$
and $\epsilon_P(\sigma)\in\Z/2$ by
\[
 \sigma(\zeta_{2n})=\zeta_{2n}^{u_\sigma},
 \qquad \sigma(\sqrt P)=(-1)^{\epsilon_P(\sigma)}\sqrt P.
\]
Let $\mathcal A\subset\Gamma$ consist of the $\sigma$
for which there are $k_P\in\Z/n$ and $e_i\in\Z/n$
satisfying
\[
 \begin{split}
 k_P&\equiv\epsilon_P(\sigma)\pmod2,\\
 (u_\sigma-1)(\delta_i+2e_i)
           +2\sum_Pw_{iP}k_P&\equiv0\pmod {2n}
                                  \quad(1\le i\le k).
 \end{split}
 \tag{A.57.3}
\]
Put $H=\langle\mathcal A\rangle\subset\Gamma$.
Then the required character group is exactly
\[
              \operatorname{Hom}(\Gamma/H,\Q/\Z).
 \tag{A.57.4}
\]
Equivalently, the largest abelian field contained
in every field factor of (A.57.1) is $B=F^H$.

Writing $n=2^r v$ with $v$ odd, one has the sharper
bound
\[
 B\subset
 \Q(\sqrt P:P\in\mathcal P)\,\Q(\mu_{2^{r+1}}).
 \tag{A.57.5}
\]
In particular every character in (A.57.4) has
order a power of two. Formula (A.57.3), or the
equivalent divisibility test below, determines
the remaining cyclotomic and quadratic relations;
complete splitting in the whole radical splitting
field imposes more conditions than (A.57.4).
\end{statement}
\begin{proof}
Put $h=n/2$, choose the positive real roots
$\alpha_P=P^{1/n}$, and put $L=F(\alpha_P:P\in\mathcal P)$.
The proof of A.47 applies to this abelian field $F$
as well: a positive rational number with an $n$th
root in $F$ is an $h$th power in $\Q$; the prime
valuations therefore give the independent Kummer
classes $(\Z/h)^{\mathcal P}$. Consequently
\[
 N=\operatorname{Gal}(L/F)\simeq(\Z/h)^{\mathcal P}.
 \tag{A.57.6}
\]
Complex conjugation fixes all the $\alpha_P$,
its commutators contain $2N$, and its commutators
with lifts of the independent quadratic characters
of $\Q(\sqrt P:P\in\mathcal P)$ generate $N/2N$.
Thus, exactly as in (A.47.3),
\[
       [\operatorname{Gal}(L/\Q),\operatorname{Gal}(L/\Q)]
                      =N,\qquad L^{\mathrm{ab}}=F.
 \tag{A.57.7}
\]
The Kummer degree assertion uses the perfect pairing
in \cite[Theorem 5.30 and Remark 5.32]{AgentMilneFT}.

Every element of $G=\operatorname{Gal}(L/\Q)$ is
uniquely specified by its restriction $\sigma$ to
$F$ and the exponents $k_P$ in
$g(\alpha_P)=\zeta_n^{k_P}\alpha_P$.
Since $\alpha_P^h=\sqrt P$, these exponents satisfy
the first line of (A.57.3). Conversely every choice
with that parity occurs: the independent degree-$h$
extensions in (A.57.6) allow all even changes of
the exponents after one lift of $\sigma$ has been
chosen. The multiplication rule is
\[
        (\sigma,k)(\tau,l)
               =(\sigma\tau,k+u_\sigma l)\pmod n.
 \tag{A.57.8}
\]
This proves the stated description as an actual
finite group, including all compatibility with
the quadratic subfields of the cyclotomic field.

The set $\Omega$ of geometric points of
$\operatorname{Spec}\mathcal E$ consists of the tuples
\[
       z_i=\zeta_{2n}^{\,\delta_i+2e_i}
                      \prod_P\alpha_P^{\,w_{iP}},
                    \qquad e_i\in\Z/n.
 \tag{A.57.9}
\]
The group action just described fixes this tuple
precisely when the second line of (A.57.3) holds.
The field factors of $\mathcal E$ are the fixed fields
of the stabilizers $G_\omega$, one for each orbit
of $G$ on $\Omega$. An abelian character trivial
on every such stabilizer is trivial on their
generated subgroup and on $[G,G]$. By (A.57.7)
it factors through $\Gamma$, and the image of
the generated stabilizers in $\Gamma$ is exactly
$H$. This proves (A.57.4) and its field formulation.
Using every point instead of one point in each
orbit changes no abelian character, since the
stabilizers in one orbit are conjugate.

For completeness, (A.57.3) is finite integer
linear algebra. Given $\sigma$ and the parity
choices $k_P$, put $d_\sigma=\gcd(n,u_\sigma-1)$.
The variables $e_i$ exist exactly when
\[
 d_\sigma\ \bigm|\ 
       \frac{u_\sigma-1}{2}\delta_i+\sum_Pw_{iP}k_P
                    \quad(1\le i\le k).
 \tag{A.57.10}
\]
Indeed, after dividing the second congruence in
(A.57.3) by two, the multiples of $u_\sigma-1$
in $\Z/n$ are exactly the multiples of $d_\sigma$.
One then computes $H$ using the multiplication
in the finite abelian group $\Gamma$.

To prove (A.57.5), put
$D=\Q(\sqrt P:P\in\mathcal P)$ and
$E_0=\Q(\alpha_P:P\in\mathcal P)$.
We have $D\subset E_0\cap F$ and
$[E_0:D]\le h^{|\mathcal P|}$ because
$\alpha_P^h=\sqrt P$.
On the other hand $E_0F=L$ has degree
$h^{|\mathcal P|}$ over $F$ by (A.57.6).
Hence $[E_0:D]=h^{|\mathcal P|}$ and $E_0$
is linearly disjoint from $F$ over $D$.
In particular
\[
 E_0\cap F=D,\qquad
 (E_0K_2)\cap F=DK_2,\quad
 K_2=\Q(\mu_{2^{r+1}})\subset F.
 \tag{A.57.11}
\]
The second equality follows by the same degree
calculation after base change from $D$ to $DK_2$.

Choose $\eta=\zeta_{2^{r+1}}$. Since $v$ is odd,
$\eta^n=-1$. The tuple
$z_i=\eta^{\delta_i}\prod_P\alpha_P^{w_{iP}}$
is one of the points in (A.57.9), and its field
is contained in $E_0K_2$. Every common abelian
subfield $B$ of the factors is contained in this
particular factor and, by (A.57.7), in $F$.
Equation (A.57.11) gives $B\subset DK_2$.
Both $D$ and $K_2$ have degree a power of two
over $\Q$, proving the assertion on character orders.
\end{proof}

\paragraph{A verified extra character and its norm identity.}
Take $n=4$, $D_1=1/5$, and $D_2=-1/5$.
In every field factor of
$\Q[Z_1,Z_2]/(Z_1^4-1/5,Z_2^4+1/5)$,
the ratio $Z_2/Z_1$ is a primitive eighth root
of unity and $Z_1^2=\pm1/\sqrt5$. Thus
every factor contains $\Q(\mu_8,\sqrt5)$.
Here this is exactly $F$, so $H=1$ and
$B=F$. The quadratic character of $\Q(\sqrt2)$
is therefore one of the required characters.
It is not generated by the two individual
squareclasses $1/5$ and $-1/5$:
their span has representatives $1,5,-1,-5$,
whereas the valuation at two of $2$ is odd.

This additional character has a constant
Brauer-class evaluation on the corresponding
old curve. More generally, for $a,b,a+b\ne0$, on
$aX^4-bY^4=1$, with $T=aX^4$ and $c=a/(a+b)$,
direct substitution gives
\[
 T-c=\frac{ab}{a+b}(X^4+Y^4),\qquad
 X^4+Y^4=(X^2+Y^2)^2-2(XY)^2.
 \tag{A.57.12}
\]
The second expression is
$N_{\Q(\sqrt2)/\Q}(X^2+Y^2+\sqrt2\,XY)$.
Consequently the quaternion class satisfies
\[
                    (T-c,2)=\left(\frac{ab}{a+b},2\right)
                         \quad\text{in }\operatorname{Br}\Q(C).
 \tag{A.57.13}
\]
For $a=3$, $b=2$ this is the fibre just described.
The constant has global reciprocity sum zero.
This calculation proves constancy for this
character by an explicit norm; it does not infer
constancy from the absence of residues.

There is also a useful exact extension of this
particular norm calculation. For $n=2^r v$, $v$ odd,
the same curve $aX^n-bY^n=1$ and the same $c$
satisfy
\[
 T-c=\frac{ab}{a+b}(X^n+Y^n),\qquad
 X^n+Y^n=
 N_{\Q(\mu_{2^{r+1}})/\Q}
                  (X^v-\zeta_{2^{r+1}}Y^v).
 \tag{A.57.14}
\]
The norm identity is the homogeneous cyclotomic
identity $\Phi_{2^{r+1}}(A/B)B^{2^r}=A^{2^r}+B^{2^r}$.
Transitivity of norms proves constant evaluation
for every character of this particular cyclotomic
field. It covers its characters of order greater
than two as well.

\paragraph{The remaining initialization step.}
For a fixed old cover
$T-c_i=a_iX_i^n$ and a new rational $c$, the
fibre algebra (A.57.1) has
$D_i=(c-c_i)/a_i$.
At a local point with $T\ne c$, put
$z=(T-c)/a_i$. Then $z+D_i=X_i^n$ and the
Steinberg relation gives the proved identity
\[
         \{T-c,D_i\}_n=\{a_i,D_i\}_n
                   \quad\text{in }K_2^M(\Q_p)/n.
 \tag{A.57.15}
\]
Indeed, apply Steinberg to $z/X_i^n$ and
$D_i/X_i^n$, whose sum is one, and expand
bilinearly. Thus each individual symbol has
the restriction of one global constant as
its local value. Equations (A.57.12)--(A.57.14)
give a further family of exact norm identities.

The calculation of all characters in (A.57.4)
does not yet prove corresponding identities
for every joint character. Nor has it produced
a new rational coefficient $a$ for which
$Y^n=(T-c)/a$ has points above the required
local points of the old cover at every place.
The fixed-finite-puncture construction A.48
does not establish this initialization with
the already chosen old powerclasses preserved.
The ordinary finite-place version of (G) with
all rational punctures therefore remains open
in this paper.

The certificate
\texttt{code/fibre\_character\_certificate.py}
compares the actual root stabilizers with
(A.57.10), computes the quotient in the
$D=(1/5,-1/5)$ example at exponents
$4,8,24,40,120$, and checks the quadratic
and full two-power cyclotomic norm identities
using exact integer determinants.
Its saved output records the finite calculations.
